\section{Reciprocity for the four corner colours}\label{sec:profiles}

The squarefree parts of the four corner contents carry a finer reciprocity
constraint than their product.  We shall prove that the product belonging
to either pair of opposite corners is congruent to $1$ or $15$ modulo
$16$.  The proof also describes the prime divisors of all six sums of two
integer distances.  Its ingredients are the Gaussian square roots of the
individual displacements and ordinary quadratic reciprocity.

\subsection{Integer distances and Gaussian roots}

Fix a strictly interior point with four positive rational distances, and use the corner
order
\[
 c_0=0,\qquad c_1=1,\qquad c_2=\mathrm i,\qquad c_3=1+\mathrm i.
\]
Write
\begin{equation}\label{eq:profile-lattice}
 x=X/L,\qquad y=Y/L,\qquad L>0,\qquad \gcd(X,Y,L)=1,
 \qquad A_j=Lr_j\in\Z_{>0}.
\end{equation}
Interchanging $x$ and $y$ if necessary, take $X$ odd.  The integral
normalization of Section~\ref{sec:integral} gives $4\mid Y$ and
$8\mid L$.  This interchange only exchanges the labels $1$ and $2$,
so it preserves the two opposite pairs as unordered pairs.

Put
\begin{equation}\label{eq:profile-colours}
 U_j=X+\mathrm iY-Lc_j,\qquad
 g_j=\gcd\bigl(|\operatorname{Re}U_j|,|\operatorname{Im}U_j|\bigr),
 \qquad m_j=\sf(g_j).
\end{equation}
The $m_j$ are positive, odd, squarefree, pairwise coprime, and coprime
to $L$.  Every prime dividing them is $1$ or $7$ modulo $8$.
These are canonical positive integers attached to the primitive
coordinate lattice.

We shall use Gaussian integers $w_j$ satisfying
\begin{equation}\label{eq:profile-integer-roots}
 U_j=m_jw_j^2,\qquad
 N_{\Q(\mathrm i)/\Q}(w_j)=\frac{A_j}{m_j}>0.
\end{equation}
For completeness, the elementary construction is as follows.  The
triple
\[
 \left(\frac{\operatorname{Re}U_j}{g_j},
       \frac{\operatorname{Im}U_j}{g_j},\frac{A_j}{g_j}\right)
\]
is a primitive Pythagorean triple whose first leg is odd and whose
second leg is even.  If its legs and hypotenuse are denoted by $a,b,k$,
the positive integers $(k+a)/2$ and $(k-a)/2$ are relatively prime,
and their product is $(b/2)^2$.  Thus they are squares, say $s^2,t^2$,
and a choice of the sign of $t$ gives $a+\mathrm ib=(s+\mathrm it)^2$.
Multiplying this root by the positive integer $\sqrt{g_j/m_j}$ gives
$w_j$.  Its norm is positive and its norm squared is $A_j^2/m_j^2$,
which proves the second identity in \eqref{eq:profile-integer-roots}.

In particular, $A_j/m_j$ is an odd sum of two integer squares, whence
\begin{equation}\label{eq:profile-radius-residue}
 A_j\equiv m_j\pmod4.
\end{equation}
The signs in the norm identity will determine the characters at primes
congruent to $7$ modulo $8$.

The diagonal pairs have squared separation $2$, and adjacent pairs have
squared separation $1$.  For $i\ne j$ define
\begin{equation}\label{eq:profile-pair-data}
 s_{ij}=m_im_j,\qquad d_{ij}=|c_i-c_j|^2\in\{1,2\},\qquad
 V_{ij}=\frac{A_i+A_j}{2}\in\Z_{>0}.
\end{equation}

\begin{lemma}[The six pair identities]\label{lem:profile-pair-identity}
For every pair $i\ne j$,
\begin{equation}\label{eq:profile-pair-identity}
 (A_i+A_j)^2-
 4s_{ij}\operatorname{Re}(w_i\overline{w_j})^2=d_{ij}L^2.
\end{equation}
Consequently, with $\Sigma_{ij}=r_i+r_j$ and
$T_{ij}=2\operatorname{Re}(w_i\overline{w_j})/L$, one has
\[
 \Sigma_{ij}^2-s_{ij}T_{ij}^2=d_{ij}.
\]
\end{lemma}

\begin{proof}
Set $Z=w_i\overline{w_j}$ and $s=s_{ij}$.  Equation
\eqref{eq:profile-integer-roots} gives
\[
 N(Z)=\frac{A_iA_j}{s},\qquad
 Z^2=\frac{U_i\overline{U_j}}{s}.
\]
The identity $2\operatorname{Re}(Z)^2=N(Z)+\operatorname{Re}(Z^2)$
therefore shows that the left side of
\eqref{eq:profile-pair-identity} is
\[
 A_i^2+A_j^2-2\operatorname{Re}(U_i\overline{U_j})
 =|U_i-U_j|^2=d_{ij}L^2.
\]
\end{proof}

\subsection{Prime divisors of all six half-sums}

\begin{theorem}[Six-pair support]\label{thm:profile-six-pair-support}
Let $p$ be an odd prime dividing $V_{ij}$.  Then
\begin{equation}\label{eq:profile-six-pair-support}
 p\nmid s_{ij},\qquad \leg{-s_{ij}}p=1.
\end{equation}
Equivalently, every odd prime divisor of $V_{ij}$ is unramified and
split in $\Q(\sqrt{-s_{ij}})$.
\end{theorem}

\begin{proof}
Write $s=s_{ij}$, $d=d_{ij}$, and $V=V_{ij}$.
If $p\mid m_i$, the point reduces to corner $i$ modulo $p$:
$U_i\equiv0$ and $A_i\equiv0$.  Since $p\nmid L$,
\[
 A_j^2\equiv dL^2\not\equiv0\pmod p.
\]
Thus $p\nmid V$.  The same argument applies to a prime dividing $m_j$,
and proves the first assertion.

Suppose next that $p\mid L$.  If $p\equiv1\pmod4$, choose a square
root of $-1$ modulo $p$ for which $X+\mathrm iY$ is nonzero.  Such a
choice exists: vanishing for both choices would force $p\mid X,Y,L$.
In this one embedding all the $U_k$ have the same nonzero residue,
and hence all $w_k$ have nonzero residues.  Reducing
\[
 \left(\frac{w_i}{w_j}\right)^2
   =\frac{m_j}{m_i}\frac{U_i}{U_j}
\]
gives $\leg{s}p=1$.  Since $\leg{-1}p=1$, the required conclusion
follows.

If $p\equiv3\pmod4$, the form $X^2+Y^2$ is anisotropic modulo $p$.
Primitivity therefore makes its value nonzero, and every $A_k$ and
$w_k$ is a unit at $p$.  Work in $\F_{p^2}=\F_p(\mathrm i)$, where
conjugation is Frobenius.  For $h=w_i/w_j$ and $a=m_j/m_i\in\F_p^*$,
we have $h^2=a$, and consequently
\[
 N_{\F_{p^2}/\F_p}(h)=h^{p+1}=a^{(p+1)/2}
     =a\leg ap.
\]
The actual norm identity in \eqref{eq:profile-integer-roots} gives at
the same time
\[
 N(h)=\frac{m_j}{m_i}\frac{A_i}{A_j}.
\]
It follows that $A_i/A_j=\leg{s}p$ modulo $p$.  The assumption
$p\mid V$ makes this ratio $-1$, so $\leg{s}p=-1$ and
$\leg{-s}p=1$.

It remains to treat $p\nmid L$.  Reducing
\eqref{eq:profile-pair-identity}, and using $p\mid V$ and $p\nmid s$,
gives
\[
 -4s\operatorname{Re}(w_i\overline{w_j})^2=dL^2.
\]
All factors on the right are nonzero modulo $p$.  For an adjacent
pair, $d=1$, this proves $\leg{-s}p=1$.  For an opposite pair the same
argument proves the assertion whenever $\leg2p=1$.

Finally, an opposite sum cannot vanish modulo a prime with
$\leg2p=-1$.  For the pair $(0,3)$, the equality $A_0=-A_3$ modulo $p$
implies $x+y=1$ in $\F_p$.  The other two distance squares then become
$2x^2$ and $2y^2$.  At least one is nonzero, which is impossible when
$2$ is a nonsquare.  For the pair $(1,2)$, the same reasoning gives
$x=y$ and the other two squares $2x^2$ and $2(x-1)^2$, again with at
least one nonzero.  This exhausts the cases.
\end{proof}

\subsection{The characters at a corner prime}

Define a character on the subgroup $\{1,7,9,15\}$ of
$(\Z/16\Z)^*$ by
\begin{equation}\label{eq:profile-sigma}
 \sigma(1)=\sigma(15)=1,\qquad
 \sigma(7)=\sigma(9)=-1.
\end{equation}
We use the same notation for positive integers supported on primes
congruent to $1$ or $7$ modulo $8$.  Multiplication of residues shows
that this is multiplicative.  For a prime in the support,
\begin{equation}\label{eq:profile-prime-sigma}
 \sigma(p)=
 \begin{cases}
 (-1)^{(p-1)/8},&p\equiv1\pmod8,\\
 (-1)^{(p+1)/8},&p\equiv7\pmod8.
 \end{cases}
\end{equation}

\begin{lemma}[Ramified character]\label{lem:profile-ramified-character}
For $p\mid m_i$ and $j\ne i$,
\begin{equation}\label{eq:profile-ramified-character}
 \leg{V_{ij}}p=\leg{m_j}p\,\sigma(p)^{d_{ij}-1}.
\end{equation}
\end{lemma}

\begin{proof}
At the corner prime $p$ we have $A_i\equiv0$ and
\[
 U_j\equiv L(c_i-c_j)\pmod p.
\]
Since $\leg2p=1$, the characters of $V_{ij}$ and $A_j$ agree.
The difference $c_i-c_j$ is a Gaussian unit for adjacent corners,
and a Gaussian unit times $1+\mathrm i$ for opposite corners.

First suppose $p\equiv7\pmod8$.  In $\F_{p^2}$, put
$u=c_i-c_j$.  The two identities for $w_j$ give
\[
 \frac{A_j}{m_j}
  =N(w_j)
  =\left(\frac{Lu}{m_j}\right)^{(p+1)/2},
\]
and hence
\begin{equation}\label{eq:profile-frobenius-radius}
 A_j=L\leg{Lm_j}p\,u^{(p+1)/2}.
\end{equation}
The exponent $(p+1)/2$ is divisible by $4$.  It therefore kills every
Gaussian unit.  In the adjacent case its final factor is $1$.
In the opposite case that factor is
\[
 z_p=(1+\mathrm i)^{(p+1)/2}
     =(-1)^{(p+1)/8}\,2^{(p+1)/4}\in\F_p,
 \qquad z_p^2=2.
\]
Because $\leg{-1}p=-1$ and $\leg2p=1$, its character is
$\leg{z_p}p=\sigma(p)$.  Taking characters in
\eqref{eq:profile-frobenius-radius}, a sign $\pm1$ has its own value
as character, so
\[
 \leg{A_j}p
 =\leg Lp\leg{Lm_j}p\,\sigma(p)^{d_{ij}-1}
 =\leg{m_j}p\,\sigma(p)^{d_{ij}-1}.
\]

Now suppose $p\equiv1\pmod8$ and work in a Gaussian embedding into
$\F_p$.  All four Gaussian units are squares.  For an adjacent pair,
$w_j^2=Lu/m_j$ therefore gives $\leg{Lm_j}p=1$, while
$A_j^2=L^2$ gives $\leg{A_j}p=\leg Lp=\leg{m_j}p$.

For an opposite pair write
$\lambda_p=\leg{1+\mathrm i}p$.  This value is unchanged on replacing
$\mathrm i$ by $-\mathrm i$, because their ratio is a Gaussian unit.
The root equation gives $\leg{Lm_j}p=\lambda_p$.  Choose
$J\in\F_p$ with $J^2=\mathrm i$.  It has order $8$, so
\[
 \leg Jp=(-1)^{(p-1)/8}=\sigma(p).
\]
Since $((1+\mathrm i)/J)^2=2$, either square root of $2$ has character
$\lambda_p\sigma(p)$; changing its sign does not change that
character.  The equation $A_j^2=2L^2$ now yields
\[
 \leg{A_j}p
  =\leg Lp\lambda_p\sigma(p)
  =\leg{m_j}p\sigma(p).
\]
This proves the formula for both congruence classes and both separations.
\end{proof}

\subsection{A pair reciprocity identity}

For a positive integer $a$, let
\begin{equation}\label{eq:profile-G}
 \mathcal G(a)=\prod_{p\equiv3\ (4)}(-1)^{v_p(a)}.
\end{equation}
Thus $\mathcal G(a)$ is $1$ or $-1$ according as the odd part of $a$
is $1$ or $3$ modulo $4$.  Let $e_j\in\{0,1\}$ be defined by
\begin{equation}\label{eq:profile-e}
 m_j\equiv(-1)^{e_j}\pmod4.
\end{equation}
Equivalently, $e_j$ is the parity of the number of prime factors of
$m_j$ which are $7$ modulo $8$.

\begin{theorem}[Pair reciprocity]\label{thm:profile-master}
For every pair $i\ne j$, let $e\in\{0,1\}$ satisfy
$e\equiv e_i+e_j\pmod2$.  Then
\begin{equation}\label{eq:profile-master}
 \mathcal G(V_{ij})^{1+e}
   =(-1)^{e_ie_j}\sigma(m_im_j)^{d_{ij}-1}.
\end{equation}
\end{theorem}

\begin{proof}
Put $s=m_im_j$ and $V=V_{ij}$.  All Jacobi symbols below have coprime
positive arguments, by Theorem~\ref{thm:profile-six-pair-support}.
Multiplication of \eqref{eq:profile-ramified-character} over the primes
of $m_i$ and $m_j$ gives
\[
 \leg Vs
  =\leg{m_j}{m_i}\leg{m_i}{m_j}\sigma(s)^{d_{ij}-1}
  =(-1)^{e_ie_j}\sigma(s)^{d_{ij}-1}.
\]
Here the second equality is quadratic reciprocity; a Jacobi symbol with
denominator $1$ is interpreted as $1$.

Let $V_{\mathrm{odd}}$ be the odd part of $V$.  The six-pair support
theorem gives, for every odd $p\mid V$,
\[
 \leg sp=\leg{-1}p.
\]
Multiplication with valuations therefore gives
\[
 \leg{s}{V_{\mathrm{odd}}}=\mathcal G(V).
\]
Also $\leg2s=1$, since every prime factor of $s$ is $1$ or $7$
modulo $8$.  A second application of quadratic reciprocity yields
\[
 \leg Vs
 =\leg{V_{\mathrm{odd}}}s
 =\leg{s}{V_{\mathrm{odd}}}\mathcal G(V)^e
 =\mathcal G(V)^{1+e}.
\]
Equating the two expressions proves the result.
\end{proof}

\subsection{The two opposite products}

\begin{theorem}[Opposite-colour congruences]\label{thm:profile-main}
For every strictly interior point with four positive rational distances,
the canonical corner colours in \eqref{eq:profile-colours} satisfy
\begin{equation}\label{eq:profile-main}
 \boxed{\quad m_0m_3\equiv1\text{ or }15\pmod{16},\qquad
               m_1m_2\equiv1\text{ or }15\pmod{16}.\quad}
\end{equation}
More precisely, an opposite product with an even number of prime factors
congruent to $7$ modulo $8$ is $1$ modulo $16$; one with an odd number
is $15$ modulo $16$.
\end{theorem}

\begin{proof}
Fix an opposite pair $i,j$ and put $s=m_im_j$ and $V=V_{ij}$.
Here $d_{ij}=2$.  If $e_i\ne e_j$, the left side of
\eqref{eq:profile-master} is $\mathcal G(V)^2=1$, and
$e_ie_j=0$.  Hence $\sigma(s)=1$.  Since $s\equiv7\pmod8$ in
this case, the conclusion is $s\equiv15\pmod{16}$.

Suppose instead that $e_i=e_j$.  By
\eqref{eq:profile-radius-residue}, the integers $A_i,A_j$ have the
same residue $(-1)^{e_i}$ modulo $4$.  Thus $V$ is odd and their
half-difference $D=(A_i-A_j)/2$ is even.  The squared-distance
identities give the exact formulas
\[
 A_0^2-A_3^2=2L(X+Y-L),\qquad
 A_2^2-A_1^2=2L(X-Y).
\]
The factors $X+Y-L$ and $X-Y$ are odd.  Consequently, in either
opposite pair,
\[
 v_2(VD)=v_2(L)-1\ge2.
\]
Since $V$ is odd, $4\mid D$.  It follows that
\[
 V\equiv A_i\equiv(-1)^{e_i}\pmod4,
 \qquad \mathcal G(V)=(-1)^{e_i}.
\]
Now $e=0$ and $e_ie_j=e_i$, so the master identity reads
\[
 (-1)^{e_i}=(-1)^{e_i}\sigma(s).
\]
Therefore $\sigma(s)=1$.  As $s\equiv1\pmod8$, this gives
$s\equiv1\pmod{16}$.
\end{proof}

Write $b=m_0m_3$, $c=m_1m_2$ and $n=bc$.  The allowable residue
patterns take the particularly simple form
\begin{equation}\label{eq:profile-residue-table}
\begin{array}{c|c}
 n\pmod{16}&(b,c)\pmod{16}\\ \hline
 1&(1,1),\ (15,15)\\
 15&(1,15),\ (15,1).
\end{array}
\end{equation}

\begin{corollary}\label{cor:profile-split}
If every prime dividing $m_0m_1m_2m_3$ is $1$ modulo $8$, then
\[
 m_0m_3\equiv m_1m_2\equiv1\pmod{16}.
\]
\end{corollary}

\begin{proof}
Every $e_j$ is zero, so both opposite products fall in the even case
of Theorem~\ref{thm:profile-main}.
\end{proof}

For adjacent pairs the exponent $d_{ij}-1$ is zero, and the same
reciprocity theorem gives
\[
 \mathcal G(V_{ij})^{1+e}=(-1)^{e_ie_j}.
\]
In particular, if two adjacent colours have the same residue modulo
$4$, the odd part of their half-sum has that same residue modulo $4$.
Together with Theorem~\ref{thm:profile-six-pair-support}, this supplies
simultaneous arithmetic information for every pair of corners.

The support conclusion also has a primitive-radius formulation.  Let
$G=\gcd(A_0,A_1,A_2,A_3)$ and write $M_j=A_j/G$.  Since all $A_j$ are
odd, $G$ and the $M_j$ are odd, and
\[
 \frac{A_i+A_j}{2}=G\frac{M_i+M_j}{2}.
\]
Every odd prime dividing $(M_i+M_j)/2$ consequently satisfies
\eqref{eq:profile-six-pair-support}.  This formulation uses the same
canonical colours \eqref{eq:profile-colours} and removes the common
content of the four integer radii.
